\section{Einleitung} \label{sec:einleitung}
\subsection{Problemstellung und Relevanz}\label{subsec:problemstellung}
Während technologische Innovationen fast alle Lebensbereiche durchdringen, hinkt die Sicherheitsinfrastruktur an vielen Schulen oft hinterher. Ein alltägliches Problem ist die Kontrolle von Schülern beim Verlassen des Schulgeländes. Bisherige Methoden erfordern meist die manuelle Überprüfung von Auslassberechtigungen durch Wachpersonal oder Sekretariat mittels Papierzetteln und Listen. Dieser Prozess ist umständlich, ineffizient und fehleranfällig. Besonders bei hohem Aufkommen nach Schulschluss oder in Freistunden stoßen analoge Werkzeuge an ihre Grenzen und verursachen unnötige Wartezeiten.

Analoge Kontrollsysteme auf Basis physischer Passierscheine sind zudem inflexibel. Ein verlorenes oder kopiertes Dokument stellt ein permanentes Sicherheitsrisiko dar, da Berechtigungen nicht digital entzogen werden können. Eine granulare, zeitbezogene Steuerung ist kaum präzise umsetzbar, zudem lassen sich physische Nachweise unbemerkt an Dritte weitergeben. Angesichts steigender Sicherheitsanforderungen und der Digitalisierung ist ein flexibles, personengebundenes und fälschungssicheres System von hoher Relevanz. Diese Arbeit fokussiert auf die Digitalisierung des Auslassprozesses, um die Zettelwirtschaft zu beenden und den Schulalltag effizienter zu gestalten, wobei das System ebenso für den regulären Zugang konzipiert ist.

\subsection{Projektbeschreibung und Ziele} \label{subsec:projektbeschreibung}
Das Projekt \enquote{Digital School Access Control} entstand im Rahmen von \enquote{Jugend Forscht 2025}, um diese Defizite durch eine moderne Lösung zu beheben. Ziel war ein Full-Stack-System, das Hardware, Software und biometrische Verifizierung vereint. Die technische Umsetzung kombiniert passive RFID-Technologie mit lokaler KI-Gesichtserkennung. Jeder Nutzer erhält eine RFID-Karte als Identifikator. Beim Einlesen wird das hinterlegte Foto auf einem Monitor für das Personal angezeigt. Ein KI-Modul vergleicht zudem das Live-Kamerabild mit dem Referenzfoto, um die Identität automatisiert zu validieren.

Das modulare System umfasst folgende Komponenten:
\begin{itemize}
    \item Ein \textbf{Node.js-Backend} mit Express-REST-API für die zentrale Logik und Datenverarbeitung.
    \item Eine \textbf{SQLite-Datenbank} für Nutzerdaten, Karten-UIDs, Berechtigungen und Zugriffsprotokolle.
    \item Eine \textbf{React-Native-App} (Expo) mit Schnittstellen für Schüler, Lehrkräfte, Wachpersonal und Administratoren.
    \item Ein \textbf{AI-Modul} auf Basis von TensorFlow.js für die Gesichtserkennung und Bildkonvertierung.
\end{itemize}
Besonderer Fokus liegt auf der Benutzerfreundlichkeit für Lehrkräfte zur intuitiven Verwaltung von Berechtigungen sowie auf einer datenschutzkonformen Umsetzung durch rein lokale Verarbeitung biometrischer Daten.

\subsection{Forschungsfrage und Hypothesen} \label{subsec:forschungsfrage}
Die zentrale Forschungsfrage lautet:
\begin{quote} \textit{Wie lässt sich die physische Zugangs- und Auslasskontrolle an Schulen durch die Kombination von passiver RFID-Technologie und lokaler KI-Gesichtserkennung effizient, fälschungssicher sowie datenschutzkonform digitalisieren, um papierbasierte Nachweisverfahren abzulösen?} \end{quote}
Zur Beantwortung wurden drei Hypothesen aufgestellt, die evaluiert werden:
\begin{itemize}
    \item \textbf{H1:} Ein kombiniertes RFID-System kann Berechtigungsentscheidungen in unter 2 Sekunden treffen und so Wartezeiten drastisch reduzieren. (Effizienz)
    \item \textbf{H2:} Lokale KI-Gesichtserkennung verhindert als zweiter Faktor die unbefugte Weitergabe von Karten, wobei eine Genauigkeit von mindestens 90\,\% erreicht wird. (Zuverlässigkeit und Fälschungssicherheit) 
    \item \textbf{H3:} Durch lokale Datenverarbeitung lassen sich zentrale DSGVO-Anforderungen (Datenminimierung, Zweckbindung) technisch unterstützen. (Datenschutz) \cite{dsgvo2016, garcia2017schoolsecurity}
\end{itemize}
